\section{Methods}
\label{sec:methods}

\subsection{Tasks, splits and protocol}
We use the four composition-only tasks of Matbench v0.1 with their official five folds, through \texttt{matbench}:
\begin{itemize}
\item steels: yield strength, 312 alloys, MAE;
\item expt\_gap: experimental band gap, 4,604 compounds, MAE;
\item expt\_is\_metal: metallicity, 4,921 compounds, classification;
\item glass: glass-forming ability, 5,680 alloys, classification.
\end{itemize}
Inside each official training fold $T_k$ the rows are split into $J=5$ inner folds (seed 0).
All selection uses only $T_k$: hyperparameters, the choice among design configurations, stack weights, the optional local correction and classification thresholds.
The design space was explored on the training rows of fold~0 only.

Two analysis plans were frozen in turn, each by SHA-256 hashes of all code, and each run scored the test folds once.
\begin{itemize}
\item \textbf{Run 1 (Amendment~1).} Originally the design was fixed on fold~0's training rows. Amendment~1 moved the choice among the explored configurations inside every training fold. An earlier, superseded run under the original plan was never scored.
\item \textbf{Run 2 (Amendment~2).} Written after run~1 had been scored.
\end{itemize}
Plans, amendments, freeze hashes, per-fold predictions and the post-hoc ablation (plan Section~12) are in the repository. The Supplementary Material (SI) gives details.

\subsection{Features}
All descriptors are deterministic functions of the formula and are computed with \texttt{matminer} \citep{ward2018matminer} and \texttt{pymatgen} \citep{ong2013pymatgen}:
\begin{itemize}
\item Magpie element-property statistics \citep{ward2016magpie} and element fractions;
\item stoichiometric, valence-orbital \citep{meredig2014}, ionic, band-centre and atomic-orbital descriptors;
\item the ``deml'', ``matminer'' and ``megnet\_el'' element-property presets;
\item Miedema enthalpies, Yang solid-solution parameters \citep{yang2012ss} and atomic packing efficiency \citep{laws2015ape}, as implemented by \citet{ward2018bmg};
\item Wen alloy descriptors and the transition-metal fraction;
\item Oliynyk element statistics \citep{oliynyk2016}.
\end{itemize}

\paragraph{DFT block (disclosed transfer learning)}
Each composition is matched by reduced formula to the Materials Project PBE data distributed with \texttt{matminer} (\texttt{matbench\_mp\_gap}, joined to \texttt{matbench\_mp\_e\_form}). The block contains:
\begin{itemize}
\item the gap, metallicity and formation energy of the lowest-formation-energy polymorph;
\item the minimum, median and maximum gap over polymorphs, the metallic fraction and the polymorph count;
\item a match indicator.
\end{itemize}
It matches 73\% of the expt\_gap and expt\_is\_metal compositions, 1\% of glass and none of steels.
The block uses no experimental labels and no pretrained model.

\subsection{Experts}
\begin{itemize}
\item \textbf{Gradient boosting.} LightGBM \citep{ke2017lightgbm} on all features, tuned inside $T_k$ by Optuna \citep{akiba2019optuna}: 20 trials, with the default configuration evaluated first. Also XGBoost \citep{chen2016xgboost} and CatBoost \citep{prokhorenkova2018catboost} with fixed settings.
\item \textbf{Tree ensembles.} Extremely randomised trees \citep{geurts2006et} on two feature sets and on element fractions alone.
\item \textbf{Kernel, neighbour and neural experts} \citep{pedregosa2011sklearn}: kernel ridge regression or RBF support-vector classifiers on element fractions and on Magpie features, distance-weighted $k$-nearest neighbours, and a two-layer perceptron.
\item \textbf{Hurdle experts (expt\_gap only).} About half of the experimental gaps are exactly zero. Each hurdle expert pairs a classifier for $E_g>0$ with a regressor fitted on the $E_g>0$ rows, and combines them as $p\,r$ or as $r\,\mathbb{1}[p\ge 0.5]$. The latter is the MAE-optimal rule for a point mass at zero plus a positive part.
\end{itemize}
CrabNet- \citep{wang2021crabnet} and MODNet-based experts were implemented but dropped, because a single fit exceeded our CPU budget.

\subsection{Stack, gated correction and thresholds}
Every expert is cross-fitted inside $T_k$ to give out-of-fold (OOF) predictions.
A simplex-constrained stack (non-negative weights summing to one) is fitted on these OOF predictions. Its loss is the MAE for regression and the Brier score for classification.

The local OOF residual correction of run~1 adds a kernel-weighted, novelty-damped average of neighbouring training residuals (SI Section~S2). In run~2 it is kept only if a nested inner-CV check confirms that it improves the honest OOF score.

Hard labels use the threshold that maximises balanced accuracy on the training OOF predictions.

\subsection{Endpoints and statistics}
The endpoints follow the board's semantics.
\begin{itemize}
\item \textbf{Regression:} MAE.
\item \textbf{Classification:} balanced accuracy, because the leaders either recorded labels (Darwin) or were scored by a \texttt{matbench} version that thresholded probabilities at 0.5 (MODNet v0.1.12), which makes their reported ``ROC-AUC'' equal to their balanced accuracy.
\item \textbf{Run 2 primary comparison:} our final pipeline versus each task's leader, paired over the same test compositions using the leaders' public per-sample predictions.
\end{itemize}

The test is a bootstrap over test points within folds ($B=10^4$), with Holm correction over the four tasks. We claim to beat a leader only if our mean is better and the Holm-adjusted $p<0.05$.
Run 1's primary endpoint (corrected versus uncorrected stack) used the same bootstrap. The same statistics are used for the post-hoc ablation, which is descriptive only.
